\section{The Certified $\Sigma$-Protocol Library}\label{schn_proof}
Every proof of the scheme comes from a library of $\Sigma$-protocols, 
verified in the Rocq theorem prover~\cite{the_coq_development_team}.
This section covers the part the scheme uses, but the artefact has many more,
including the Okamoto, linear-relation, and the
Parallel, And, Eq, and Neq compositions.

In an implementation, group elements and field elements are both large integers and
mixing them up is a known source of bugs~\cite{10.1007/978-3-662-63958-0_24}. 
The library separates
them by type: an abstract group $G$, an abstract field $F$, and a
\rocq[\emph{vector space}]{vector_space} structure relating them, with scalar multiplication
written $\mbox{\textasciicircum} : G \to F \to G$ because it is
exponentiation in a multiplicative group. Every protocol and
theorem is generic in this structure, so the development instantiates to
Schnorr groups and elliptic curves alike, and the field and group laws
automate the proofs. Listing~\ref{schn_nat}
(Appendix~\ref{appendix_listings}) shows the Schnorr protocol (\rocq{schnorr_protocol}, \rocq{schnorr_simulator}, \rocq{accepting_conversation}, \rocq{special_honest_verifier_zkp_enum}): the prover's transcript for randomness $u$
and witness $x$, the simulator's transcript for randomness $u$ without
the witness, and the verifier's test $g^r = a \cdot h^c$.


Completeness and special soundness are plain equations over the group
and field, with no probability involved (\rocq{schnorr_completeness}
and \rocq{special_soundness_berry}). Special soundness extracts
the witness $(r - r')/(c - c')$ from two accepting transcripts that share
an announcement but have different challenges. Zero-knowledge
compares two probability distributions and so needs a model of them. Following~\cite{erwig2006functional}, a \rocq[distribution]{dist} is
a list of (value, probability) pairs, and two distributions are equal when
they are permutations of each other. The type is a monad, so a
distribution over $n$ independent uniform draws is the $n$-fold
composition of the \rocq[uniform distribution]{uniform_with_replacement}. The \rocq[Schnorr distribution]{schnorr_distribution} binds a
uniform $u$ and returns the honest transcript, and the \rocq[simulator distribution]{simulator_distribution}
binds a uniform $u$ and returns the simulated one. The theorem of
Listing~\ref{schn_nat} says that, when \texttt{lf} enumerates the field,
the two are permutations of each other, i.e.\ equal as distributions. The
proof is the textbook one ($u \mapsto u + c\cdot x$ is a bijection of
$\mathbb{Z}_q$). The statement is as strong as the informal
definition, because the transcripts themselves coincide, and a simulator that
ignores its input cannot satisfy it. Every composition in the library has a
zero-knowledge theorem of this shape, proved by a bijection on the
randomness, coordinatewise for conjunctions and for Chaum--Pedersen, and
by a reshuffling of the simulated sub-challenges for disjunctions.

\textbf{The three compositions the scheme uses:}
\rocqfile[OrSigmaGen]{src/Crypto/OrSigmaGen.v} is the $n$-ary
disjunction of Cramer, Damg\aa rd, and Schoenmakers~\cite{cramer1994proofs}
over Schnorr statements $h_i = g_i^{x_i}$. The prover, knowing the
witness of branch $k$, commits honestly there, simulates every other
branch with chosen sub-challenges, and sets $c_k = c - \sum_{i\ne k}
c_i$. The verifier checks $c = \sum_i c_i$ and every branch equation (soundness \rocq{generalised_or_sigma_soundness_main}, zero-knowledge \rocq{generalised_or_special_honest_verifier_zkp_perm}).
\rocqfile[EncProof]{src/Crypto/EncProof.v} is the membership proof for ElGamal: a ciphertext
$(g^r, g^v h^r)$ encrypts one of the messages $v_1,\ldots,v_n$, as an
$n$-way disjunction of Chaum--Pedersen statements, one per message (\rocq{generalised_construct_encryption_proof_elgamal_real}, \rocq{generalised_accepting_encryption_proof_elgamal}).
\rocqfile[DecProof]{src/Crypto/DecProof.v} is the Chaum--Pedersen proof that a ciphertext decrypts
to a given plaintext under the key $h = g^x$ (\rocq{construct_decryption_proof_elgamal_real}, \rocq{decryption_proof_accepting_conversations}). Each comes with completeness,
special soundness, and zero-knowledge theorems. The statements
of \rocqfile[OrSigmaGen]{src/Crypto/OrSigmaGen.v} are over one group. The Schulze certificate applies
it to the product group $G \times G$, treating a ciphertext as one
element, which the vector-space abstraction allows without a new proof
(Listing~\ref{or_proof}, Appendix~\ref{appendix_listings}).

\textbf{Instantiation:}
To run the extracted code we instantiate the vector space. The \rocq[Schnorr
group]{Schnorr_group} is a record of an integer $v$ with proofs $0 < v < p$ and $v^q = 1
\bmod p$, the field $\mathbb{Z}_q$ a \rocq[record]{Zp} of a reduced integer, and
\rocq{pow} the scalar multiplication between them (\rocq{pow_vspace})
(Listings~\ref{group_def}--\ref{vec_space}, Appendix~\ref{appendix_listings}).
The group and field laws are proved with the primality of $p$ and $q$ as
section hypotheses, which Coqprime discharges for the 2048-bit Helios
parameters~\cite{adida2008helios} with primality
certificates~\cite{thery2007primality}. Mixing a group element with a
field element is a type error.

